% draft_abstract.tex
\begin{abstract}
This paper estimates the causal effect of daily Twitter sentiment on intraday stock returns
using two identification strategies: firm fixed-effects OLS and Double Machine Learning (DoubleML)
with XGBoost nuisance learning. The primary dataset is a Bloomberg daily panel of 503 S\&P~500
firms spanning January 2025 to March 2026 (approximately 160{,}000 firm-day observations),
extended to a 1.19-million-observation S\&P~500 panel and a 1.97-million-observation Russell 3000 panel
over 2016--2026 for temporal and cross-sectional analysis. In the S\&P~500 sample, Twitter sentiment
is statistically insignificant in all fixed-effects specifications; the DoubleML estimate is
$-0.924$ ($p < 0.01$), but a placebo test---regressing yesterday's return on today's
sentiment---yields a coefficient of $+2.348$ ($p < 0.01$), indicating substantial reverse
causality from prices to sentiment. News sentiment is the more robust predictor across both
frameworks ($-0.759^{***}$ in DoubleML). A Fama-MacBeth replication of Gu and Kurov (2020)
produces insignificant coefficients on lagged Twitter sentiment in every calendar year from
2016 to 2025 for the S\&P~500, but significant effects in the broader Russell 3000.
Cap-group decomposition reveals a sharp gradient: the Panel OLS sentiment coefficient is
$+0.375^{***}$ for large-cap firms, rising to $+0.804^{***}$ for mid-cap and
$+0.942^{***}$ for small-cap firms, consistent with an information-environment explanation
in which Twitter provides marginal predictive content primarily where institutional analyst
coverage is sparse. Daily long-short portfolios sorted on sentiment generate positive Sharpe
ratios in most years, suggesting the signal is concentrated in the tails of the
cross-sectional distribution rather than in a linear average slope.
\end{abstract}
